\begin{table}[htbp]\centering
\caption{Classification accuracy under increasing sampling noise
(\(n=3\), \(\varepsilon=0.5\), averaged over 30 realisations). Test points
are perturbed by i.i.d.\ noise \(\sigma\cdot\mathcal{N}(0,I)\) on the
coordinate torus. Since \(\Lambda\) and \(\Theta\) are functions of the known
metric, they are exactly invariant under this perturbation in exact
arithmetic; only \(\mathrm{Dgm}_1\) changes with \(\sigma\). At
\(\varepsilon=0.5\) the metric deformation dominates the sampling noise for
all \(\sigma\le 0.20\), and all four descriptors remain at \(100\%\). This
table therefore serves as a control experiment confirming that the
classification pipeline is robust to noise at moderate deformation
magnitudes.}
\label{tab:noise-robustness}
\begin{tabular}{c|cccc}\hline
\(\sigma\) & Lambda & Theta & Diagram & Joint \\\hline
0.00 & 1.00 (0.00) & 1.00 (0.00) & 1.00 (0.00) & 1.00 (0.00) \\
0.05 & 1.00 (0.00) & 1.00 (0.00) & 1.00 (0.00) & 1.00 (0.00) \\
0.10 & 1.00 (0.00) & 1.00 (0.00) & 1.00 (0.00) & 1.00 (0.00) \\
0.15 & 1.00 (0.00) & 1.00 (0.00) & 1.00 (0.00) & 1.00 (0.00) \\
0.20 & 1.00 (0.00) & 1.00 (0.00) & 1.00 (0.00) & 1.00 (0.00) \\
\hline\end{tabular}\end{table}